\chapter{Aprendizado de Máquina Causal}

\section{Da previsão aos efeitos de tratamento}

Estimadores de aprendizado de máquina também podem ser usados para
inferência causal quando o objetivo é estimar um efeito de tratamento
$Y(1) - Y(0)$. Os comandos deste capítulo visam duas quantidades:

\begin{itemize}
\item o \textbf{Efeito Médio do Tratamento} (ATE)
\item o \textbf{Efeito Médio do Tratamento Condicional} (CATE),
      $\tau(x) = \E[Y(1) - Y(0) \mid X = x]$
\end{itemize}

Todos exigem um indicador de tratamento \texttt{d} e um conjunto de covariáveis
\texttt{x}. Eles não substituem um desenho de pesquisa claro, mas podem ajudar
a descobrir efeitos heterogêneos e melhorar a precisão.

\section{Causal Forest}

Causal Forest (Athey e Wager, 2018) adapta Random Forest para estimar
efeitos de tratamento heterogêneos. Divide as covariáveis para
maximizar a variância dos efeitos estimados entre folhas.

\begin{lstlisting}[caption={Causal Forest}]
set_seed(42)
let n = 1000
generate df x1 = rnormal(n)
generate df x2 = rnormal(n)
generate df d = (0.5 + 0.3*x1 > runiform(n))
generate df y = 1 + 2*x1 - x2 + d*(0.5 + 0.4*x1) + rnormal(n)

let m_cf = causalforest(y ~ d, df, x="x1,x2", trees=200, depth=5)
print(m_cf)
\end{lstlisting}

A saída inclui:

\begin{itemize}
\item o ATE estimado
\item uma distribuição do CATE (média, quartis)
\item importância das variáveis
\end{itemize}

\section{DR-Learner e TMLE}

O \textbf{DR-Learner} (Kennedy, 2023) e o \textbf{TMLE} (van der Laan e
Rubin, 2006) são abordagens duplamente robustas. Cada um constrói um
pseudo-resultado a partir de funções nuisance estimadas e depois
regredide o pseudo-resultado sobre as covariáveis para recuperar o CATE
ou ATE.

\begin{lstlisting}[caption={DR-Learner}]
let m_dr = dr_learner(y ~ d, df, x="x1,x2", folds=5)
print(m_dr)
\end{lstlisting}

\begin{lstlisting}[caption={TMLE}]
let m_tmle = tmle(y ~ d, df, w="x1,x2")
print(m_tmle)
\end{lstlisting}

Métodos duplamente robustos são atraentes porque o ATE permanece
consistente se o escore de propensão ou o modelo de resultado estiver
bem especificado.

\section{Orthogonal Random Forest}

ORF (Oprescu, Syrgkanis e Wu, 2019) residualisa o tratamento e o
resultado em relação às confundidoras antes de rodar uma floresta nos
sinais residualisados.

\begin{lstlisting}[caption={Orthogonal Random Forest}]
let m_orf = orf(y ~ d, df, x="x1,x2", w="x1,x2", trees=100, depth=5)
print(m_orf)
\end{lstlisting}

É útil quando o tratamento não é randomizado e a heterogeneidade
relevante vem de um subconjunto de covariáveis.

\section{Controle Sintético Bayesiano}

\texttt{bsc} implementa uma versão bayesiana do estimador de controle
sintético. É apropriado quando uma única unidade tratada é comparada a
várias unidades doadoras e o objetivo é um contrafactual com
intervalos credíveis.

\begin{lstlisting}[caption={Controle Sintético Bayesiano}]
// df tem colunas y, c1, c2, ..., t
let m_bsc = bsc(df, y, "c1,c2", 25, prior=1.0)
print(m_bsc)
\end{lstlisting}

A saída reporta o efeito do tratamento $\tau$, seu desvio padrão, um
intervalo credível de 95\% e os pesos do controle sintético.

\section{Double Machine Learning}

\texttt{dml\_crossfit} estima o ATE no modelo parcialmente linear

\[
Y = \theta D + g(X) + \varepsilon
\]

usando validação cruzada K-fold para a função nuisance $g(X)$.

\begin{lstlisting}[caption={Double Machine Learning}]
let m_dml = dml_crossfit(y ~ d, df, x="x1,x2", folds=5)
print(m_dml)
\end{lstlisting}

Como a nuisance é estimada em um fold separado, o ATE resultante tem
distribuição de taxa $\sqrt{n}$ padrão e erros padrão honestos.

\section{Escolhendo um método}

\begin{itemize}
\item \textbf{Efeitos heterogêneos:} Causal Forest ou ORF
\item \textbf{ATE duplamente robusto:} DR-Learner, TMLE, DML
\item \textbf{Unidade única tratada:} BSC
\end{itemize}

Todos esses comandos são \textit{experimentais} e exigem que o Hayashi
seja compilado com \texttt{--features experimental}. A validação contra
implementações de referência é parcial ou não suportada para alguns.
Use-os para exploração e como complementos a um desenho claro, não como
substitutos.
